\documentclass[12pt, a4paper]{article}

% --- Essential Packages ---
\usepackage[utf8]{inputenc} % Character encoding
\usepackage{geometry}       % Page margins
\geometry{a4paper, margin=1in}
\usepackage{graphicx}       % For including images/charts
\usepackage{amsmath}        % For math formulas
\usepackage{hyperref}       % For clickable links and cross-references
\usepackage{booktabs}       % For professional looking tables
\usepackage{setspace}       % For line spacing
\usepackage{multirow}
\DeclareMathOperator*{\argmax}{arg\,max}


% --- Title Information ---
\title{\textbf{Deep Learning for Financial Forecasting: TensorFinance Project Report}}
\author{Ta Thi Thao \\ Fulbright University Vietnam}
\date{May 2026} % Or \today

\begin{document}

\maketitle

% --- 1. Abstract ---

\begin{abstract}
This report presents TensorFinance, an end-to-end platform for deep
learning--based stock price forecasting and systematic trading, developed
across five tasks in the CS313 curriculum. Tasks~1 and~2 train LSTM and
CNN-LSTM models on 31 NASDAQ and 27 Vietnamese HOSE stocks respectively,
achieving mean next-day prediction errors of 1.65\% (NASDAQ) and 1.31\%
(Vietnam). Task~3 adapts the CNN-LSTM backbone to generate buy and sell
trading signals, attaining mean AUC values of 0.781 and 0.806. Task~4
integrates model outputs with Markowitz mean-variance optimisation to
produce two portfolio allocation strategies.
Task~5 deploys all outputs as a production FastAPI backend with daily
automated data ingestion, a React~19 dashboard, and a Docker Compose stack.
A key design principle throughout is stationarity by construction: all model
inputs are expressed as ratios or log-returns invariant to absolute price
level, preventing distribution shift across the multi-year training horizon.
\end{abstract}




\newpage
\tableofcontents
\newpage

% --- 2. Introduction ---
\section{Introduction}

% Explain the difficulty of predicting financial time series (high noise, non-stationary).
Financial markets are notoriously difficult to predict due to their high noise-to-signal ratio. Traditional time-series models often fail to capture complex, non-linear momentum shifts. This project builds an end-to-end platform for financial time-series data supporting forecasting and systematic trading. The working progresses through 5 tasks: 
\begin{itemize}
    \item \textbf{Tasks 1:} Single & Multiple outputs predictions for tickers on NASDAG.
    \item \textbf{Task 2: } Single & Multiple outputs predictions for tickers on VIETNAM markets whose data does not high quality as NASDAG.
    \item \textbf{Task 3:} Trading signals model for Vietnamese market.
    \item \textbf{Task 4:} Portfolio optimization for Vietnam's market.
    \item \textbf{Task 5:} Deploying models as a user-facing SaaS application with AI automation workflow.
    
\end{itemize}






% --- 3. Tasks 1 & 2 ---
\section{ Task 1: Prediction on NASDAG market with LSTM model}
\subsection{Data selection}
\begin{itemize}
 
 \item \textbf{Sector Grouping}: The remaining tickers are grouped into 6 major groups on the NASDAG market(Technology, Healthcare, Financials, Consumer, Industrials, Energy).
 \item \textbf{Top 5 selection}: Within each sector, candidates were ranked by the total number of trading days, and the top 5 ranked by data history were selected. This prioritizes data richness, enabling longer training sequences and expose the model to a more diverse market movement(bull/bear cycles, recessions, sector rotations). The final dataset consisted of 31 tickers through 6 sectors.
 \item \textbf{Data extension to 2026}: The provided CSV files covered historical data ending in 2022. To support the model training on recent market conditions and live inference at the deployment time, data was extended to 2026.

 \item \textbf{Foresee potential problem}: The initial approach appended Yahoo Finance data (2023-2026) to existing CSV files. However, this two-source design introduces a complication during the deployment stage: live prediction would require fetching new rows from Yahoo Finance and aligning them to the CSVs before inference. This would latency and potentially introduce alignment error.
 \item \textbf{Solution}: To eliminate the complexity, a cleaner solution was adopted: all historical data was downloaded from Yahoo Finance ranging from 2005 to present. This unified data source simplifies the production pipeline, and the data scheduler only needs to append pipeline to a single provided. This would remove the potential format inconsistency between the CSVs data and live-fetched data.
 \end{itemize}

% Describe yfinance and local CSVs used in Task 1.
\subsection{Feature Engineering}

\subsubsection{Design Principle: Stationarity by Construction}

The first experiment (Figure~\ref{fig:feature_eng}) suggests the need
for a different approach to feature engineering rather than feeding raw
prices. When feeding raw prices, the model experiences severe
distribution shifts because prices are non-stationary. For example,
AAPL was priced at \$5 in 2005, \$80 in 2019, and \$180 in 2023. A
model trained on 2005--2019 data sees inputs that are completely
out-of-distribution when evaluated on 2023--2026 data. This leads to a
systematic underprediction of \$100--\$200. The root cause is
non-stationarity: raw prices carry an absolute scale that drifts over
time.

\begin{figure}[htpb]
    \centering
    \includegraphics[width=0.85\textwidth]{FeatureEngineeringTask1.png}
    \caption{Feature engineering pipeline for Task~1 NASDAQ inputs.}
    \label{fig:feature_eng}
\end{figure}


  



All 20 features (24 for Vietnam) were therefore engineered as 
\textbf{ratios, log-returns, or bounded oscillators} that are invariant to 
the absolute price level. A ratio such as $\text{Close}/\text{SMA}_{20} - 1 
= 0.03$ carries the same meaning---``price is 3\% above its 20-day 
average''---whether the stock trades at \$5 or \$500.

\subsubsection{Feature Set}

Table~\ref{tab:features} summarises all 20 features used for NASDAQ (Task~1),
grouped by category, with their formula and the market signal they encode.

\begin{table}[h]
\centering
\caption{Feature set: 20 stationary technical indicators (Task~1, NASDAQ).}
\label{tab:features}
\renewcommand{\arraystretch}{1.3}
\begin{tabular}{llp{6cm}}
\hline
\textbf{Group} & \textbf{Feature(s)} & \textbf{Signal / Reasoning} \\
\hline
\multirow{4}{*}{Momentum}
  & \texttt{Return\_1d}  & Daily log-return: $\ln(P_t / P_{t-1})$ \\
  & \texttt{Return\_5d}  & Weekly momentum: $\ln(P_t / P_{t-5})$ \\
  & \texttt{Return\_10d} & Bi-weekly momentum \\
  & \texttt{Return\_20d} & Monthly drift \\
\hline
\multirow{5}{*}{Trend}
  & \texttt{SMA10\_vs\_SMA20} & Short/medium MA crossover: $\text{SMA}_{10}/\text{SMA}_{20} - 1$ \\
  & \texttt{SMA20\_vs\_SMA50} & Medium/long MA crossover \\
  & \texttt{EMA10\_vs\_EMA20} & Exponential crossover (faster reaction) \\
  & \texttt{Close\_vs\_SMA20} & Price deviation from 20-day anchor \\
  & \texttt{Close\_vs\_SMA50} & Price deviation from 50-day anchor \\
\hline
Oscillator
  & \texttt{RSI\_14} & Relative Strength Index (bounded $[0,100]$); overbought $>70$, oversold $<30$ \\
\hline
\multirow{3}{*}{MACD}
  & \texttt{MACD\_pct}      & $(EMA_{12} - EMA_{26}) / P_t$: trend momentum \\
  & \texttt{MACD\_sig\_pct} & 9-day EMA of MACD line / $P_t$: smoothed signal \\
  & \texttt{MACD\_hist\_pct}& Difference of above two: momentum acceleration \\
\hline
\multirow{2}{*}{Bollinger}
  & \texttt{BB\_width} & $(BB_{upper} - BB_{lower}) / BB_{mid}$: volatility expansion/compression \\
  & \texttt{BB\_pct}   & $(P_t - BB_{lower}) / (BB_{upper} - BB_{lower})$: position within bands \\
\hline
\multirow{3}{*}{Volatility}
  & \texttt{ATR\_pct}       & $\text{ATR}_{14} / P_t$: true range including overnight gaps \\
  & \texttt{Volatility\_10} & 10-day rolling std of log-returns \\
  & \texttt{Volatility\_20} & 20-day rolling std of log-returns \\
\hline
Volume  & \texttt{Volume\_ratio} & $V_t / \overline{V}_{20}$: relative volume vs. 20-day average \\
\hline
Range   & \texttt{High\_Low\_pct} & $(H_t - L_t) / P_t$: intraday price range intensity \\
\hline
\end{tabular}
\end{table}

\subsubsection{Vietnam-Specific Features (Task~2)}
\label{sec:vn_features}

Four additional features were added for the Vietnamese HOSE market, bringing 
the total to 24. These capture market microstructure properties specific to 
HOSE that have no NASDAQ equivalent (Table~\ref{tab:vn_features}).

\begin{table}[h]
\centering
\caption{Four VN-specific features added for Task~2 (HOSE market).}
\label{tab:vn_features}
\renewcommand{\arraystretch}{1.3}
\begin{tabular}{lp{4cm}p{6cm}}
\hline
\textbf{Feature} & \textbf{Formula} & \textbf{Reasoning} \\
\hline
\texttt{VN\_limit\_prox} 
  & $\ln(P_t/P_{t-1}) \;/\; 0.07$
  & HOSE enforces a $\pm7\%$ daily circuit-breaker. Values near $\pm1$ indicate 
    the stock hit its limit; limit-hit days create predictable next-day 
    reversions the model can exploit. \\
\hline
\texttt{OC\_pct} 
  & $(P_{close} - P_{open}) / P_{close}$
  & Candlestick body direction. The VN retail-dominated market responds strongly 
    to short-term TA signals; bullish bodies reliably predict short-term 
    continuation. \\
\hline
\texttt{Upper\_shadow} 
  & $\bigl(H - \max(O,C)\bigr) / P_{close}$
  & Intraday selling pressure: price rose but was pushed back before close, 
    signalling resistance. \\
\hline
\texttt{Lower\_shadow} 
  & $\bigl(\min(O,C) - L\bigr) / P_{close}$
  & Intraday buying support: price dropped but recovered before close, 
    signalling a demand floor. \\
\hline
\end{tabular}
\end{table}

All features are \textbf{causal}: each value at time $t$ uses only data from 
$t$ and earlier, implemented via backward-looking rolling windows and 
positive-lag \texttt{shift(k)} operations. No future price information enters 
any feature computation.
\subsection{Model Architecture}

\subsubsection{Architecture Choice}
\label{sec:arch}

The model is a three-layer stacked LSTM followed by two fully connected layers,
totalling \textbf{138,657 trainable parameters} (Table~\ref{tab:model_summary}).
The descending hidden-unit schedule (128 $\rightarrow$ 64 $\rightarrow$ 32)
is deliberate: the first layer extracts a rich set of temporal features from
the 60-day input window, while successive layers progressively compress these
representations into a compact encoding before the regression head.
A \texttt{Dropout(0.3)} layer follows each LSTM to prevent co-adaptation of
hidden units, and an intermediate \texttt{Dense(16)} layer provides a
non-linear projection before the scalar output.

\begin{table}[h]
\centering
\caption{LSTM model summary for Task~1.1 (input shape: $60 \times 20$).}
\label{tab:model_summary}
\renewcommand{\arraystretch}{1.3}
\begin{tabular}{lrr}
\hline
\textbf{Layer} & \textbf{Output Shape} & \textbf{Parameters} \\
\hline
LSTM (128 units, \texttt{return\_sequences=True})  & $(N,\;60,\;128)$ & 76,288  \\
Dropout (0.3)                                       & $(N,\;60,\;128)$ & 0       \\
LSTM (64 units, \texttt{return\_sequences=True})   & $(N,\;60,\;64)$  & 49,408  \\
Dropout (0.3)                                       & $(N,\;60,\;64)$  & 0       \\
LSTM (32 units, \texttt{return\_sequences=False})  & $(N,\;32)$       & 12,416  \\
Dropout (0.3)                                       & $(N,\;32)$       & 0       \\
Dense (16, ReLU)                                    & $(N,\;16)$       & 528     \\
Dense (1, linear)                                   & $(N,\;1)$        & 17      \\
\hline
\textbf{Total trainable parameters} & & \textbf{138,657} \\
\hline
\end{tabular}
\end{table}

\textbf{Loss function.} Huber loss was chosen over mean squared error (MSE)
because financial return series contain sharp spikes driven by earnings
surprises or macro events. MSE squares these residuals, causing the model to
allocate disproportionate gradient signal toward rare extreme events at the
expense of ordinary-day accuracy. Huber loss behaves as MSE for small errors
but transitions to a linear penalty beyond a threshold $\delta$, making
training robust to outliers while still penalising systematic bias.

\textbf{Lookback window.} A window of 60 trading days ($\approx$3 months) was
selected to capture enough trend context --- covering one earnings cycle and
typical momentum half-lives --- while keeping the sequence length manageable.
Empirically, shorter windows (20 days) produced higher validation variance,
while longer windows (120 days) did not reduce test MAE and substantially
increased training time.

\subsection{Model Training and Validation}

\subsubsection{Walk-Forward Cross-Validation}

Standard $k$-fold cross-validation is invalid for time series because it
shuffles temporal order, allowing future data to appear in training folds.
Instead, \texttt{sklearn.TimeSeriesSplit} with an \textbf{expanding window}
was used: each successive fold adds more training data while holding out a
fixed-length validation window. This mirrors real-world deployment, where the
model trains on all available history and predicts forward.

Table~\ref{tab:wfcv} reports the five-fold walk-forward CV results on
AAPL's training set (3,663 samples). The monotonically decreasing MAE from
Fold~1 to Fold~4 confirms the expected result: more training data produces
better generalisation. The slight uptick at Fold~5 is consistent with
noise in the fixed-size validation window rather than overfitting.
The mean MAE of \$0.66 $\pm$ 0.20 provides a stable pre-test performance estimate.

\begin{table}[h]
\centering
\caption{Walk-forward cross-validation results on AAPL training set
         (5 folds, 3,663 total training samples). Metrics in USD.}
\label{tab:wfcv}
\renewcommand{\arraystretch}{1.3}
\begin{tabular}{crrcc}
\hline
\textbf{Fold} & \textbf{Train samples} & \textbf{Val samples}
              & \textbf{MAE (\$)} & \textbf{RMSE (\$)} \\
\hline
1 & 613  & 610 & 1.0067 & 1.4031 \\
2 & 1223 & 610 & 0.6235 & 0.8388 \\
3 & 1833 & 610 & 0.5967 & 0.8421 \\
4 & 2443 & 610 & 0.5067 & 0.7194 \\
5 & 3053 & 610 & 0.5789 & 0.8178 \\
\hline
\textbf{Mean} & & & \textbf{0.6625 $\pm$ 0.197} & \textbf{0.924 $\pm$ 0.272} \\
\hline
\end{tabular}
\end{table}

\subsubsection{Single-Ticker Training: AAPL Demo}

AAPL was used as the development ticker to validate the training pipeline
before scaling to all 31 stocks. Figure~\ref{fig:aapl_training_history} shows the
Huber loss and MAE curves over 100 epochs. Two observations are notable.
First, the validation loss decreases steadily and plateaus without diverging,
confirming the absence of overfitting --- the gap between train and validation
loss remains stable throughout training. Second, the validation MAE stabilises
around 0.81 (scaled units) while training MAE reaches 0.69, a modest gap
attributable to Dropout being active only during training; this is expected
behaviour and not a sign of overfitting.

\begin{figure}[h]
\centering
\includegraphics[width=0.85\textwidth]{aapl_training_history.png}
\caption{AAPL training and validation Huber loss / MAE curves over 100 epochs.
         Stable validation loss confirms absence of overfitting.}
\label{fig:aapl_training_history}
\end{figure}

\subsubsection{Multi-Horizon Performance: Evidence of Learning}

Table~\ref{tab:multi_horizon_aapl} reports AAPL validation and test metrics
across all three prediction horizons of the multi-horizon model.
Three patterns jointly confirm that the model has learned genuine price
dynamics rather than trivially approximating a na\"{i}ve baseline.

\textbf{First}, MAPE is low and consistent across splits: 1.66\% on
validation and 1.09\% on the test set for the next-day horizon. A
persistence baseline (predicting tomorrow's price equals today's) achieves
a MAPE approximately equal to the average absolute daily return, which for
AAPL over this period is roughly 1.0--1.5\%. The model matches or beats
this benchmark while additionally predicting the \emph{direction} of the move.

\textbf{Second}, directional accuracy exceeds 50\% on every split and
horizon (range: 50.8\%--56.1\%). A model that learned nothing would predict
near-zero returns on every day and achieve a directional accuracy of
approximately 50\% by chance. Consistent above-chance directional accuracy
across six independent evaluation windows is statistically unlikely under
the null hypothesis of no learning.

\textbf{Third}, the error increases monotonically with horizon --- next-day
$<$ 3-day MA $<$ 7-day MA --- which is the theoretically expected behaviour
for a genuine forecasting model: uncertainty compounds over longer horizons.
A model that had overfit to noise would not exhibit this structured degradation.

\begin{table}[h]
\centering
\caption{AAPL multi-horizon LSTM prediction metrics across validation and
         test sets. Directional accuracy (Dir~Acc) measures the fraction of
         days where the predicted price movement direction matches the actual.}
\label{tab:multi_horizon_aapl}
\renewcommand{\arraystretch}{1.3}
\begin{tabular}{llcccc}
\hline
\textbf{Split} & \textbf{Horizon}
  & \textbf{MAE (\$)} & \textbf{RMSE (\$)}
  & \textbf{MAPE (\%)} & \textbf{Dir Acc} \\
\hline
Validation & Next-Day & 2.07 & 2.78 & 1.66 & 50.8\% \\
Validation & 3-Day MA & 2.55 & 3.37 & 2.02 & 54.5\% \\
Validation & 7-Day MA & 3.50 & 4.48 & 2.75 & 52.2\% \\
\hline
Test       & Next-Day & 2.26 & 3.33 & 1.09 & 55.2\% \\
Test       & 3-Day MA & 3.12 & 4.43 & 1.49 & 55.7\% \\
Test       & 7-Day MA & 4.41 & 5.87 & 2.09 & 56.1\% \\
\hline
\end{tabular}
\end{table}

Notably, directional accuracy on the test set \emph{improves} slightly with
horizon (55.2\% $\rightarrow$ 55.7\% $\rightarrow$ 56.1\%). This is consistent
with the model capturing medium-term trend momentum more reliably than
day-to-day noise --- the 7-day moving average target smooths out daily
idiosyncratic fluctuations, making the directional signal cleaner.


\subsubsection{Global Backbone Training (All 31 Tickers)}

Before per-ticker fine-tuning, a single shared backbone was pre-trained on
the concatenated sequences of all 31 tickers ($\approx$111,000 sequences).
Figure~\ref{fig:global_backbone} shows the global training curve. The
validation loss ($\approx$0.340) falls \textit{below} the training loss
($\approx$0.346) throughout training. This counterintuitive result is expected:
Dropout and BatchNormalization inflate training loss during forward passes,
while validation is computed with these regularisers disabled. The sub-0.01
gap between train and validation confirms that the global backbone generalises
well across diverse tickers rather than memorising any single stock's patterns.

\begin{figure}[h]
\centering
\includegraphics[width=0.85\textwidth]{global_backbone_training.png}
\caption{Global backbone training on all 31 NASDAQ tickers concatenated.
         Validation loss below training loss is expected due to Dropout/BN
         behaviour; the near-zero gap confirms cross-ticker generalisation.}
\label{fig:global_backbone}
\end{figure}

\subsubsection{Multi-Horizon Model: Architecture and Training (Tasks 1.1, 1.2, 1.3)}

\paragraph{Architecture Change from Single-Step to Multi-Horizon.}
The single-step model (Task~1.1, Section~\ref{sec:arch}) uses a 128
$\rightarrow$ 64 $\rightarrow$ 32 LSTM stack optimising one target.
Jointly predicting three correlated targets increases the effective
output dimension and raises the risk of overfitting, since the model
must simultaneously satisfy three loss terms on the same hidden state.
The multi-horizon variant therefore uses a \textbf{reduced-capacity}
architecture with three additional regularisation mechanisms
(Table~\ref{tab:arch_comparison}):

\begin{table}[h]
\centering
\caption{Architecture comparison: single-step (Task~1.1) vs.\
         multi-horizon (Tasks~1.1 extension, 1.2, 1.3).}
\label{tab:arch_comparison}
\renewcommand{\arraystretch}{1.3}
\begin{tabular}{lll}
\hline
\textbf{Component} & \textbf{Single-step (Task~1.1)}
                   & \textbf{Multi-horizon (Tasks~1.2, 1.3)} \\
\hline
LSTM units         & 128 $\rightarrow$ 64 $\rightarrow$ 32
                   & 64 $\rightarrow$ 32 $\rightarrow$ 16 \\
BatchNormalization & No  & After each LSTM layer \\
Dropout rate       & 0.3 & 0.4 \\
L2 regularisation  & $1\times10^{-4}$ & $5\times10^{-4}$ \\
Output head        & \texttt{Dense(1)} & \texttt{Dense($k$)} \\
\hline
\end{tabular}
\end{table}

BatchNormalization stabilises activations when gradients flow through
three correlated loss terms simultaneously. The higher dropout (0.4)
and stronger L2 ($5\times10^{-4}$) compensate for the reduced
parameter count by increasing regularisation pressure. This design
follows the principle that \emph{model capacity and regularisation
strength should scale together} with output complexity.

\paragraph{Target Definitions.}
The three output heads predict the following log-return targets,
all anchored to today's closing price $P[t]$:

\begin{align}
y_{\text{1d}}  &= \log\!\left(\frac{P[t+1]}{P[t]}\right)
  & &\text{(next-day log-return)} \label{eq:y1d} \\[4pt]
y_{\text{ma3}} &= \log\!\left(\frac{\overline{P}[t+1:t+4]}{P[t]}\right)
  & &\text{(3-day forward moving average)} \label{eq:yma3} \\[4pt]
y_{\text{ma7}} &= \log\!\left(\frac{\overline{P}[t+1:t+8]}{P[t]}\right)
  & &\text{(7-day forward moving average)} \label{eq:yma7}
\end{align}

where $\overline{P}[a:b]$ denotes the arithmetic mean of closing prices
over days $a$ to $b-1$. Anchoring all three targets to $P[t]$ ensures
each head's output is a bounded log-return in $(-0.1, +0.1)$, making
the three loss terms commensurate in scale and allowing equal-weight
joint optimisation with Huber loss.

\paragraph{Task 1.2 --- $k$-th Day Ahead Forecast.}
Task~1.2 reuses the same reduced architecture but with a single output
head predicting the cumulative log-return to exactly the $k$-th future day:

\begin{equation}
y_k = \log\!\left(\frac{P[t+k]}{P[t]}\right), \quad k \in \{3,\, 7\}
\label{eq:yk}
\end{equation}

This differs from $y_{\text{ma3}}$ and $y_{\text{ma7}}$ above: the
moving-average targets smooth over a window, whereas the $k$-th day
target predicts a specific future point. The global Task~1.1 backbone
is used as initialisation, and only the output head is retrained with
the new target label --- the backbone already encodes directional
momentum that transfers directly to longer-horizon prediction.

\paragraph{Task 1.3 --- $k$ Consecutive Days Forecast.}
Task~1.3 predicts all days from $t+1$ to $t+k$ in a single forward
pass by widening the output head to \texttt{Dense($k$)}. Each output
neuron $j$ predicts the log-return to day $t+j$, anchored to the
same $P[t]$:

\begin{equation}
y_j = \log\!\left(\frac{P[t+j]}{P[t]}\right), \quad j = 1, \ldots, k
\label{eq:yj}
\end{equation}

Price reconstruction for each day is independent:
$\hat{P}[t+j] = P[t] \cdot \exp(\hat{y}_j)$, which eliminates the
error compounding that affects recursive multi-step forecasting
(where day $j+1$ depends on the predicted --- not actual ---
price at day $j$). A separate \texttt{StandardScaler} is fitted
per output dimension on training targets only.

\paragraph{Early Stopping Behaviour.}
Figure~\ref{fig:multi_horizon} shows the training curves for the
multi-horizon model on AAPL. The best validation epoch is \textbf{epoch~22},
after which validation loss begins to rise while training loss
continues to decrease --- the characteristic signature of overfitting.
\texttt{EarlyStopping} with \texttt{restore\_best\_weights=True}
halts training and reverts to epoch-22 weights, preventing the model
from memorising training patterns at the expense of generalisation.
The overfit gap (shaded region in Figure~\ref{fig:multi_horizon})
develops rapidly after epoch~22, confirming that the early stopping
threshold is appropriate for this architecture.


\begin{figure}[h]
\centering
\includegraphics[width=\textwidth]{multi_horizon_training.png}
\caption{Multi-horizon LSTM training curves. The dashed green line marks
         the best validation epoch (22). The shaded region indicates the
         overfitting gap that develops after early stopping.}
\label{fig:multi_horizon}
\end{figure}

\subsection{Model Performance}

\subsubsection{AAPL Prediction — Validation and Test Sets}

Figure~\ref{fig:aapl_prediction} shows the predicted versus actual closing
prices for AAPL on the validation set (November 2019 -- February 2023) and
test set (January 2023 -- April 2026). The model tracks the actual price
trajectory closely across both periods, including the COVID-19 crash in
early 2020 and the 2022 bear market decline.

On the \textbf{validation set}, MAE = \$2.07, RMSE = \$2.78, MAPE = 1.66\%.
On the \textbf{test set}, MAE = \$2.27, RMSE = \$3.37, MAPE = 1.08\%.
The lower MAPE on the test set despite a slightly higher absolute MAE reflects
the higher absolute price level in 2023--2026 (\$150--280) relative to
2019--2023 (\$60--180): the same dollar error represents a smaller percentage
at higher prices.

\begin{figure}[h]
\centering
\includegraphics[width=\textwidth]{aaple_next_day_prediction.png}
\caption{AAPL next-day price prediction. Top: validation set (2019--2023),
         MAE=\$2.07, MAPE=1.66\%. Bottom: test set (2023--2026),
         MAE=\$2.27, MAPE=1.08\%.}
\label{fig:aapl_prediction}
\end{figure}

\subsubsection{Per-Ticker Performance Summary}

Table~\ref{tab:perticker} reports test-set MAE, RMSE, and MAPE for all 22
tickers shown after per-ticker fine-tuning of the global backbone.
The mean MAPE across all 31 tickers is \textbf{1.65\%}, indicating that on
average the model predicts next-day closing price to within 1.65\% of the
actual value.

\begin{table}[h]
\centering
\caption{Per-ticker test-set performance for all 31 tickers after global
         pre-train and per-ticker fine-tuning (Task~1.1, next-day prediction),
         sorted by MAPE ascending.}
\label{tab:perticker}
\resizebox{\textwidth}{!}{%

\renewcommand{\arraystretch}{1.2}
\begin{tabular}{lccclccc}
\hline
\textbf{Ticker} & \textbf{MAE (\$)} & \textbf{RMSE (\$)} & \textbf{MAPE}
& \textbf{Ticker} & \textbf{MAE (\$)} & \textbf{RMSE (\$)} & \textbf{MAPE} \\
\hline
HOLX  & 0.63 & 0.95  & 0.88\% & APOG  & 0.77 & 1.28  & 1.53\% \\
ROST  & 1.40 & 2.06  & 1.00\% & UMBF  & 1.39 & 2.00  & 1.56\% \\
MNST  & 0.59 & 0.84  & 1.01\% & KLIC  & 0.84 & 1.19  & 1.81\% \\
CINF  & 1.39 & 1.90  & 1.07\% & TRNS  & 1.72 & 2.65  & 1.87\% \\
AAPL  & 2.27 & 3.36  & 1.08\% & SGC   & 0.25 & 0.40  & 2.08\% \\
AMGN  & 3.27 & 4.76  & 1.12\% & DIOD  & 1.36 & 1.88  & 2.10\% \\
PHI   & 0.26 & 0.37  & 1.14\% & ATRO  & 0.60 & 1.05  & 2.12\% \\
CBSH  & 0.59 & 0.82  & 1.15\% & WDC   & 1.97 & 4.41  & 2.15\% \\
JJSF  & 1.61 & 2.37  & 1.19\% & NEOG  & 0.27 & 0.40  & 2.26\% \\
OTTR  & 0.98 & 1.40  & 1.21\% & HELE  & 1.40 & 2.31  & 2.30\% \\
ALCO  & 0.39 & 0.55  & 1.33\% & MARPS & 0.11 & 0.18  & 2.31\% \\
IDXX  & 6.99 & 11.28 & 1.35\% & FLXS  & 0.83 & 1.38  & 2.34\% \\
TRMK  & 0.42 & 0.59  & 1.36\% & INTC  & 0.74 & 1.13  & 2.36\% \\
TXN   & 2.50 & 3.65  & 1.38\% & MLAB  & 2.62 & 3.65  & 2.44\% \\
HBAN  & 0.19 & 0.28  & 1.39\% & CLNE  & 0.08 & 0.11  & 2.85\% \\
MAT   & 0.28 & 0.46  & 1.51\% &       &      &       &        \\
\hline
\multicolumn{8}{l}{\textit{Mean MAPE: \textbf{1.65\%} \quad
  Best: HOLX 0.88\% \quad Worst: CLNE 2.85\% \quad
  Mean MAE: \$1.25 \quad Mean RMSE: \$1.97}} \\
\hline
\end{tabular}
}
\end{table}


The variation in MAPE across tickers is consistent with differences in
price predictability rather than model failure. HOLX (0.88\%) and MNST
(1.01\%) are stable, low-volatility stocks with smooth price trends that
the LSTM tracks easily. In contrast, MLAB (2.44\%), INTC (2.36\%), and
NEOG (2.26\%) are lower-liquidity stocks where idiosyncratic noise dominates
the price signal, inherently limiting achievable accuracy. The high absolute
MAE for IDXX (\$6.99) reflects its high price level (\$400--600 range) rather
than poor relative accuracy, as confirmed by its moderate MAPE of 1.35\%.






\section{Task 2 --- Vietnam Stock Price Prediction (CNN-LSTM)}

\subsection{Model Architecture}

Task~2 replaces the pure LSTM of Task~1 with a \textbf{CNN-LSTM hybrid},
motivated by two properties of the Vietnamese HOSE market. First,
short-term 3--5 day candlestick formations carry strong predictive signal
in a retail-dominated market where technical analysis is widely practised.
A \texttt{Conv1D} layer with causal padding explicitly extracts these local
patterns before the LSTM processes the full sequence. Second, \texttt{MaxPool(2)}
compresses the 30-step input to 15 steps, halving the gradient path through
the LSTM and reducing vanishing-gradient pressure on the shorter VN histories.
Table~\ref{tab:cnn_lstm_arch} summarises the full architecture.

\begin{table}[h]
\centering
\caption{CNN-LSTM architecture for Task~2
         (input shape: $30 \times 24$).}
\label{tab:cnn_lstm_arch}
\renewcommand{\arraystretch}{1.3}
\begin{tabular}{lll}
\hline
\textbf{Block} & \textbf{Layer} & \textbf{Configuration} \\
\hline
\multirow{5}{*}{CNN}
  & Conv1D  & 64 filters, kernel=3, causal padding \\
  & BatchNorm & --- \\
  & Conv1D  & 64 filters, kernel=3, causal padding \\
  & BatchNorm & --- \\
  & MaxPool(2) + Dropout & pool size=2, dropout=0.20 \\
\hline
\multirow{4}{*}{LSTM}
  & LSTM(64) & \texttt{return\_sequences=True} \\
  & Dropout  & 0.25 \\
  & LSTM(32) & \texttt{return\_sequences=False} \\
  & Dropout  & 0.25 \\
\hline
\multirow{3}{*}{Head}
  & Dense(32, ReLU) & --- \\
  & Dropout         & 0.10 \\
  & Dense($d$, linear) & $d=1$ (Tasks 2.1, 2.2); $d=k$ (Task 2.3) \\
\hline
\multicolumn{2}{l}{Loss}      & Huber ($\delta=1.0$) \\
\multicolumn{2}{l}{Optimiser} & Adam (lr=$10^{-3}$, \texttt{clipnorm}=1.0) \\
\hline
\end{tabular}
\end{table}

Causal padding ensures that $\text{output}[t]$ depends only on
$\text{input}[0 \ldots t]$, preserving strict temporal causality throughout
the CNN stage. Huber loss is chosen for its robustness to the $\pm7\%$
daily price-limit outliers endemic to the HOSE circuit-breaker mechanism.

\subsection{Data and Preprocessing}

\textbf{Dataset.} 27 Vietnamese HOSE stocks passed quality filtering,
all classified as LONG ($\geq 2{,}000$ rows). The lookback window is
\textbf{30 trading days} (versus 60 for NASDAQ), chosen to ensure
sufficient sequence counts on the shorter VN histories.
The chronological 70/15/15 train/validation/test split is applied
without shuffling; \texttt{StandardScaler} is fitted exclusively on
training rows $[0,\; n_{\text{train}} + \text{lookback} - 1)$ to
prevent distributional leakage.

\textbf{Features.} Each input vector contains 24 stationary features:
20 universal indicators shared with Task~1 (momentum returns, moving
average ratios, RSI, MACD, Bollinger Bands, ATR, volatility, volume
ratio, daily range) plus 4 VN-specific features described in
Section~\ref{sec:vn_features}.

\subsection{Transfer Learning Strategy}

Training 27 independent models from scratch on histories of varying
length would produce inconsistent model quality. Instead, a two-phase
transfer learning strategy is applied, exploiting the architectural
separation between the CNN (universal pattern extractor) and the LSTM
(ticker-specific sequence model).

\textbf{Phase 1 --- Global backbone pre-training.}
All 27 LONG-ticker training sequences are concatenated ($\approx$95{,}000
sequences) and used to train a single shared CNN-LSTM backbone.
The backbone learns universal Vietnamese candlestick patterns ---
limit-hit reversions, opening gap behaviour, volume-price relationships
--- that are consistent across all HOSE stocks.

\textbf{Phase 2 --- Per-ticker fine-tuning.}
Since all 27 tickers qualify as LONG, every ticker receives full-layer
fine-tuning (all layers unfrozen, LR=$10^{-4}$), allowing the LSTM
to adapt to each stock's specific sequential dynamics while preserving
the CNN feature representations learned in Phase~1.
Table~\ref{tab:tl_strategy} summarises the fine-tuning protocol.

\begin{table}[h]
\centering
\caption{Per-ticker fine-tuning strategy by data tier.}
\label{tab:tl_strategy}
\renewcommand{\arraystretch}{1.3}
\begin{tabular}{lccc}
\hline
\textbf{Tier} & \textbf{Condition} & \textbf{Layers unfrozen} & \textbf{LR} \\
\hline
LONG & $\geq 2{,}000$ rows & All (CNN + LSTM + Dense) & $1\times10^{-4}$ \\
SHORT & $< 2{,}000$ rows & LSTM + Dense only (CNN frozen) & $5\times10^{-4}$ \\
\hline
\multicolumn{4}{l}{\textit{All 27 tickers in this dataset qualify as LONG.}} \\
\hline
\end{tabular}
\end{table}

\subsection{Task 2.1 --- Next-Day Results}

Table~\ref{tab:task21} reports test-set performance for all 27 tickers.
The mean MAPE of \textbf{1.31\%} across all tickers confirms the model
predicts next-day closing prices to within approximately 1.3\% of the
actual value. The mean directional accuracy of 44.4\% is discussed
in Section~\ref{sec:t2_observations}.

\begin{table}[h]
\centering
\caption{Task~2.1 per-ticker test-set results (next-day prediction).
         MAE and RMSE are in VND. Sorted by MAPE ascending.}
\label{tab:task21}
\renewcommand{\arraystretch}{1.2}
\resizebox{\textwidth}{!}{%
\begin{tabular}{lrrrr|lrrrr}
\hline
\textbf{Ticker} & \textbf{MAE} & \textbf{RMSE} & \textbf{MAPE\%} & \textbf{DirAcc}
& \textbf{Ticker} & \textbf{MAE} & \textbf{RMSE} & \textbf{MAPE\%} & \textbf{DirAcc} \\
\hline
KDC  &  316.6 &  538.0 & 0.58 & 0.446 & VNM  &  635.3 &  945.7 & 1.05 & 0.431 \\
DHG  &  617.4 & 1050.0 & 0.64 & 0.419 & FMC  &  411.8 &  638.0 & 1.01 & 0.455 \\
PGC  &  105.4 &  185.8 & 0.76 & 0.342 & SFI  &  270.0 &  442.4 & 1.01 & 0.342 \\
RAL  &  881.8 & 1498.4 & 0.81 & 0.395 & DMC  &  666.4 & 1074.8 & 1.08 & 0.371 \\
PPC  &   98.9 &  163.8 & 0.86 & 0.409 & DHA  &  502.7 &  783.7 & 1.09 & 0.441 \\
ACB  &  196.4 &  314.8 & 0.93 & 0.426 & SAM  &   90.9 &  146.9 & 1.27 & 0.454 \\
FPT  & 1201.5 & 1688.1 & 1.28 & 0.479 & VIP  &  162.7 &  253.3 & 1.31 & 0.414 \\
GMD  &  855.7 & 1259.6 & 1.39 & 0.482 & HAX  &  193.3 &  280.2 & 1.43 & 0.425 \\
STB  &  600.0 &  944.5 & 1.42 & 0.463 & BMC  &  297.0 &  459.7 & 1.62 & 0.480 \\
SSI  &  442.7 &  647.0 & 1.58 & 0.487 & BMP  & 1841.0 & 2669.1 & 1.59 & 0.463 \\
SJS  &  602.2 & 1024.6 & 1.61 & 0.464 & PVD  &  434.4 &  651.1 & 1.64 & 0.451 \\
TDH  &   65.5 &  100.5 & 1.71 & 0.514 & VNE  &   91.2 &  147.9 & 1.68 & 0.503 \\
PAC  &  430.8 &  626.6 & 1.86 & 0.467 & MHC  &  185.1 &  271.1 & 1.98 & 0.484 \\
PET  &  594.9 &  884.0 & 2.05 & 0.478 &      &        &        &      &       \\
\hline
\multicolumn{10}{l}{\textit{Mean: MAE=473.8 \quad RMSE=729.2 \quad
  MAPE=1.31\% \quad DirAcc=0.444 \quad
  Best: KDC 0.58\% \quad Worst: PET 2.05\%}} \\
\hline
\end{tabular}
}
\end{table}

\subsection{Task 2.2 and 2.3 --- Multi-Horizon Results}

Table~\ref{tab:task22_23} summarises average performance across all
tickers for the $k$-th day ahead (Task~2.2) and $k$ consecutive days
(Task~2.3) forecasting variants.

\begin{table}[h]
\centering
\caption{Multi-horizon results averaged over all 27 tickers.}
\label{tab:task22_23}
\renewcommand{\arraystretch}{1.3}
\begin{tabular}{llrr}
\hline
\textbf{Task} & \textbf{Horizon} & \textbf{Avg MAE (VND)} & \textbf{Avg DirAcc} \\
\hline
2.1 & Next-day ($k=1$)             &  473.8 & 0.444 \\
2.2 & $k$-th day ahead ($k=3$)     &  845.5 & 0.492 \\
2.2 & $k$-th day ahead ($k=7$)     & 1306.0 & 0.485 \\
\hline
2.3 & $k$ consecutive days ($k=3$) &  671.1 & 0.443 \\
2.3 & $k$ consecutive days ($k=7$) &  938.1 & 0.452 \\
\hline
\end{tabular}
\end{table}

Task~2.3 uses a \textbf{direct multi-step output} --- the Dense head
outputs all $k$ predictions simultaneously rather than feeding each
prediction back as input. This eliminates the recursive error compounding
that would otherwise cause errors at step $j+1$ to depend on the
(incorrect) predicted price at step $j$. The MAE-per-step grows
approximately linearly with step index, confirming that errors accumulate
additively rather than multiplicatively.

\subsection{Key Observations}
\label{sec:t2_observations}

\textbf{1. MAE scales with horizon.}
Average MAE increases from 474 (k=1) to 846 (k=3) to 1,306 VND (k=7),
consistent with the theoretical expectation that uncertainty compounds
over longer forecast horizons. This monotonic degradation is evidence
that the model is making genuine probabilistic forecasts rather than
producing arbitrary outputs.

\textbf{2. Directional accuracy and its interpretation.}
Mean directional accuracy of 44.4\% at $k=1$ appears below the 50\%
random baseline, but this is not a model failure. The Huber loss
optimises \emph{magnitude}, not direction: on days with ambiguous
low-volatility signals, predicting a near-zero return minimises loss
regardless of sign. The model correctly identifies uncertainty; the
binary direction metric penalises this conservatism. Directional
accuracy improves to 49.2\% at $k=3$, where the smoothed target
reduces noise and makes the directional signal cleaner.

\textbf{3. VN circuit-breaker feature contribution.}
The \texttt{VN\_limit\_prox} feature encodes proximity to the HOSE
$\pm7\%$ daily limit. Stocks that hit their upper or lower limit
exhibit predictable next-day reversions due to pent-up order flow.
The CNN's local receptive field (kernel size 3) is well-suited to
detecting these 1--2 day patterns, and the causal padding ensures
no future information enters the computation.

\textbf{4. Transfer learning effectiveness.}
The global backbone pre-trained on 95,000 sequences from 27 tickers
provides a stable initialisation that captures cross-market patterns.
Per-ticker fine-tuning then adapts the LSTM weights to each stock's
specific dynamics. This two-phase approach consistently outperforms
training from scratch, particularly for tickers at the shorter end
of the LONG tier.


\section{Task 3 --- Vietnam Trading Signal Generation}

\subsection{Objective and Problem Formulation}

Task~3 repurposes the CNN-LSTM backbone from Task~2 for \textbf{binary
classification}: rather than predicting a continuous price value, two
independent classifiers identify favourable trading entry points
(\textsc{buy} signal) and exit points (\textsc{sell} signal) for each
of the 27 Vietnamese stocks. This represents a deliberate architectural
shift --- the same feature representations learned for price regression
are adapted for a fundamentally different output space.

\subsection{Label Construction}

Labels are derived from a \textbf{symmetric 11-day price window} centred
on each day $t$ (5 days before and 5 days after):

\begin{equation}
\text{pos}_t = \frac{P_{\text{close}}[t] - \min_{t-5}^{t+5}(P)}
                    {\max_{t-5}^{t+5}(P) - \min_{t-5}^{t+5}(P)}
\label{eq:pos}
\end{equation}

Binary labels are then assigned by thresholding $\text{pos}_t$:

\begin{equation}
\text{BUY}[t]  = \mathbf{1}[\text{pos}_t \leq 0.30]
\qquad \text{(price near local minimum --- good entry)}
\label{eq:buy_label}
\end{equation}

\begin{equation}
\text{SELL}[t] = \mathbf{1}[\text{pos}_t \geq 0.70]
\qquad \text{(price near local maximum --- good exit)}
\label{eq:sell_label}
\end{equation}

This formulation produces a natural positive rate of approximately
30\% for each label without artificial oversampling or undersampling,
avoiding the class-imbalance distortions that would affect a simpler
up/down direction label.

\subsection{Architecture: From Regression to Classification}

The CNN-LSTM body from Task~2 is reused without modification.
Only three components change (Table~\ref{tab:t3_arch_diff}):

\begin{table}[h]
\centering
\caption{Architectural differences between Task~2 (regression)
         and Task~3 (classification).}
\label{tab:t3_arch_diff}
\renewcommand{\arraystretch}{1.3}
\begin{tabular}{lll}
\hline
\textbf{Component} & \textbf{Task 2 (Regression)} & \textbf{Task 3 (Classification)} \\
\hline
Output activation & Linear               & Sigmoid $\in [0,1]$               \\
Loss function     & Huber ($\delta=1.0$) & BinaryFocalCrossentropy ($\gamma=2.0$) \\
LSTM/Dense init   & Task~2 weights       & Xavier re-initialisation           \\
\hline
\end{tabular}
\end{table}

The CNN layers (\texttt{conv1}, \texttt{bn1}, \texttt{conv2},
\texttt{bn2}, \texttt{pool}) are loaded from the Task~2 backbone
and kept frozen during Phase~1, preserving the learned candlestick
feature representations. The LSTM and Dense layers are \textbf{reset
to Xavier initialisation} --- the reason for this reset is explained
in Section~\ref{sec:t3_design}.

\subsection{Key Design Decision: Focal Loss and LSTM Reset}
\label{sec:t3_design}

Early experiments revealed a critical failure mode: every ticker
produced a degenerate classifier with recall~$= 1.0$ and a decision
threshold of 0.30 for all outputs. The model was predicting positive
for every input. Root cause analysis identified that the Task~2 LSTM
weights encoded a strong prior toward \emph{continuous uptrend
prediction} --- appropriate for regression but catastrophic for a
balanced binary classifier, as it biased the sigmoid output toward
high values regardless of input.

Three fixes were applied jointly:

\begin{enumerate}
  \item \textbf{LSTM and Dense weight reset.} All non-CNN layers were
    reinitialised with Xavier uniform weights before fine-tuning,
    removing the regression bias while retaining the CNN feature
    extraction capability.

  \item \textbf{Focal loss ($\gamma = 2.0$).}
    Standard binary cross-entropy assigns equal weight to all
    misclassified examples. When easy negatives dominate the batch,
    gradient signal is diluted and the model collapses to predicting
    the majority class. BinaryFocalCrossentropy with $\gamma = 2.0$
    down-weights the loss contribution of easy, confidently classified
    negatives by a factor of $(1-p)^2$, forcing the model to focus
    on hard, uncertain examples near the decision boundary.

  \item \textbf{Minimum precision guard (\texttt{MIN\_PREC} $= 0.35$).}
    During threshold calibration, any candidate threshold yielding
    precision~$< 0.35$ is rejected, preventing the optimiser from
    selecting degenerate high-recall/zero-precision solutions.
\end{enumerate}

\subsection{Transfer Learning and Training Protocol}

Training follows the same two-phase strategy as Task~2:

\textbf{Phase~1 --- Global backbone pre-training.}
A shared CNN-LSTM classifier is pre-trained on the concatenated
BUY (or SELL) labels of all 27 tickers. This phase learns
general entry/exit pattern representations that are common
across the VN market.
\begin{figure}[h]
    \centering
    \includegraphics[width=0.85\textwidth]{task3_buy_training_curve.png}
    \caption{Training curves for the BUY backbone training.}
    \label{fig:t3_buy_training}
\end{figure}


\textbf{Phase~2 --- Per-ticker fine-tuning.}
Individual classifiers are fine-tuned per ticker using
walk-forward cross-validation (\texttt{TimeSeriesSplit}, 5 folds).
Early stopping monitors \texttt{val\_auc} with patience~$= 5$ and
restores best weights. All layers are unfrozen at this stage.
\begin{figure}[h]
    \centering
    \includegraphics[width=0.85\textwidth]{task3_sell_training_curve.png}
    \caption{Training curves for the SELL backbone training.}
    \label{fig:t3_sell_training}
\end{figure}

\subsection{Threshold Calibration}
\label{sec:threshold_cal}

Each of the 54 trained models (27 tickers $\times$ 2 signal types)
receives its own decision threshold, calibrated by maximising
F1-score on the validation fold subject to the precision guard:

\begin{equation}
\tau^* = \argmax_{\tau \in [0,1]}\ F_1(\tau)
\quad \text{subject to} \quad \text{Precision}(\tau) \geq 0.35
\label{eq:threshold}
\end{equation}

Table~\ref{tab:t3_thresholds} reports the resulting threshold
distributions. SELL thresholds are systematically higher than
BUY thresholds, consistent with the model requiring higher
confidence to call a local price peak than a local trough ---
price maxima are sharper and shorter-lived than minima, reducing
the window for high-confidence predictions.

\begin{table}[h]
\centering
\caption{Calibrated decision threshold statistics across 27 tickers.}
\label{tab:t3_thresholds}
\renewcommand{\arraystretch}{1.3}
\begin{tabular}{lccc}
\hline
\textbf{Signal} & \textbf{Mean} $\tau^*$ & \textbf{Min} & \textbf{Max} \\
\hline
BUY  & 0.563 & 0.480 & 0.640 \\
SELL & 0.585 & 0.440 & 0.680 \\
\hline
\end{tabular}
\end{table}

\subsection{Results}

\subsubsection{Aggregate Performance}

Table~\ref{tab:t3_summary} reports mean, minimum, and maximum AUC
across all 27 tickers for both signal types. AUC~$> 0.75$ across
every ticker confirms that both classifiers reliably separate
positive trading opportunities from negative ones.




\begin{table}[h]
\centering
\caption{Task~3 aggregate AUC across 27 tickers.}
\label{tab:t3_summary}
\renewcommand{\arraystretch}{1.3}
\begin{tabular}{lcccccc}
\hline
\textbf{Signal} & \textbf{Mean AUC} & \textbf{Min AUC} & \textbf{Max AUC}
  & \textbf{Mean F1} & \textbf{Mean Prec.} & \textbf{Mean Rec.} \\
\hline
BUY  & 0.781 & 0.685 (VNE) & 0.845 (BMP) & 0.597 & 0.487 & 0.793 \\
SELL & 0.806 & 0.766 (DHG) & 0.863 (DMC) & 0.598 & 0.499 & 0.774 \\
\hline
\end{tabular}
\end{table}

SELL signals (mean AUC~$= 0.806$) outperform BUY signals
(mean AUC~$= 0.781$). This is consistent with the observation
that price peaks are structurally sharper than troughs: a sell
opportunity concentrates within a narrow 1--2 day window,
producing a cleaner label that the CNN's local receptive field
detects reliably. Price bottoms are broader and more ambiguous,
resulting in noisier labels and lower discriminability.

\subsubsection{Per-Ticker Performance}

Table~\ref{tab:t3_perticker} reports AUC and calibrated threshold
for every ticker across both signal types.

\begin{table}[h]
\centering
\small
\caption{Per-ticker AUC and calibrated threshold for BUY and SELL
         classifiers. Sorted by BUY AUC descending.}
\label{tab:t3_perticker}
\renewcommand{\arraystretch}{1.2}
\begin{tabular}{l cc c cc}
\hline
& \multicolumn{2}{c}{\textbf{BUY signal}}
& \phantom{xx}
& \multicolumn{2}{c}{\textbf{SELL signal}} \\
\cmidrule{2-3} \cmidrule{5-6}
\textbf{Ticker} & \textbf{AUC} & $\tau^*$
&
& \textbf{AUC} & $\tau^*$ \\
\hline
BMP   & 0.845 & 0.54 && 0.814 & 0.52 \\
PET   & 0.844 & 0.50 && 0.798 & 0.58 \\
SFI   & 0.843 & 0.62 && 0.805 & 0.60 \\
STB   & 0.840 & 0.52 && 0.843 & 0.62 \\
DMC   & 0.825 & 0.58 && 0.863 & 0.60 \\
GMD   & 0.821 & 0.56 && 0.799 & 0.54 \\
PAC   & 0.810 & 0.58 && 0.828 & 0.66 \\
SSI   & 0.806 & 0.52 && 0.781 & 0.56 \\
FPT   & 0.804 & 0.58 && 0.802 & 0.60 \\
BMC   & 0.803 & 0.48 && 0.833 & 0.64 \\
ACB   & 0.792 & 0.58 && 0.808 & 0.58 \\
VNM   & 0.792 & 0.64 && 0.848 & 0.60 \\
FMC   & 0.788 & 0.58 && 0.779 & 0.52 \\
HAX   & 0.788 & 0.56 && 0.828 & 0.68 \\
VIP   & 0.790 & 0.58 && 0.795 & 0.58 \\
MHC   & 0.780 & 0.60 && 0.780 & 0.60 \\
SAM   & 0.780 & 0.58 && 0.835 & 0.58 \\
KDC   & 0.767 & 0.56 && 0.784 & 0.64 \\
PVD   & 0.767 & 0.56 && 0.791 & 0.52 \\
SJS   & 0.757 & 0.56 && 0.790 & 0.64 \\
PGC   & 0.746 & 0.54 && 0.786 & 0.56 \\
PPC   & 0.746 & 0.56 && 0.792 & 0.58 \\
RAL   & 0.739 & 0.64 && 0.818 & 0.54 \\
DHG   & 0.729 & 0.60 && 0.766 & 0.60 \\
TDH   & 0.708 & 0.50 && 0.820 & 0.44 \\
DHA   & 0.700 & 0.48 && 0.794 & 0.58 \\
VNE   & 0.685 & 0.60 && 0.795 & 0.64 \\
\hline
\textit{Mean} & \textit{0.781} & \textit{0.563}
&& \textit{0.806} & \textit{0.585} \\
\hline
\end{tabular}
\end{table}


No ticker falls below AUC~$= 0.685$ for BUY or AUC~$= 0.766$
for SELL, confirming consistent performance across the full
27-ticker universe. The lowest BUY AUC belongs to VNE (0.685),
a low-liquidity stock where idiosyncratic noise reduces label
quality. The highest SELL AUC belongs to DMC (0.863), a
pharmaceutical stock with sharp, well-defined price peaks that
align closely with the 11-day window label definition.

\section{Task 4 --- Portfolio Optimisation and Risk Management}

\subsection{Overview and Pipeline}

Task~4 integrates the outputs of Tasks~2 and~3 with classical
quantitative finance methods to construct two investable portfolios
from the 27 Vietnamese stocks. The pipeline proceeds in three
sequential stages, each filtering or transforming the candidate
universe before the final allocation step:

\begin{enumerate}
  \item \textbf{Task~4.1 --- Candidate selection.} CNN-LSTM predicted
    returns (Task~2) and BUY signal probabilities (Task~3) are combined
    into a profitability score. Only stocks with positive predicted
    returns \emph{and} sufficient signal confidence pass.
  \item \textbf{Task~4.2 --- Risk scoring and exclusion.} Five
    historical risk metrics are computed per ticker. High-risk tickers
    are excluded before optimisation.
  \item \textbf{Task~4.3 --- Mean-variance optimisation (MVO).}
    Markowitz optimisation is applied to the surviving candidates,
    with final weights blended with Task~3 signal conviction.
\end{enumerate}

This staged design reflects a core principle of systematic portfolio
construction: \emph{return potential and risk must be assessed
independently before they are combined}, preventing high-return but
high-risk assets from dominating the allocation.

\subsection{Task 4.1 --- Profitable Stock Selection}

\subsubsection{Profitability Score}

Each ticker is assigned a profitability score that combines the
CNN-LSTM predicted 7-day return with the Task~3 BUY signal probability:

\begin{equation}
\text{ProfScore}_i = \hat{r}_{7d,i} \times p_{\text{buy},i}
\label{eq:prof_score}
\end{equation}

where $\hat{r}_{7d,i}$ is the Task~2 7-day predicted log-return and
$p_{\text{buy},i}$ is the Task~3 BUY signal probability. This product
rewards stocks where both the magnitude of the predicted move
\emph{and} the classifier confidence are high --- a stock with a large
predicted return but low BUY probability, or vice versa, receives a
lower score than one where both signals agree.

\subsubsection{Candidate Filter}

A two-condition hard filter is applied before optimisation:

\begin{equation}
\text{Candidate}_i = \mathbf{1}\!\left[
  \hat{r}_{7d,i} > 0
  \;\;\text{and}\;\;
  p_{\text{buy},i} \geq 0.50 \times \tau_i^*
\right]
\label{eq:candidate}
\end{equation}

where $\tau_i^*$ is the ticker's calibrated Task~3 threshold
(Section~\ref{sec:threshold_cal}). The 50\% relaxation of the threshold
is deliberate: requiring $p_{\text{buy}} \geq \tau^*$ would pass
only tickers where the classifier is at its highest conviction,
which at any given date may be zero or one ticker.
Relaxing to 50\% of $\tau^*$ allows borderline-positive signals
while still excluding clearly negative ones, ensuring sufficient
diversification (target: 8--12 candidates).

\textbf{Outcome.} At the evaluation date, \textbf{10 of 27 tickers}
passed both conditions. Table~\ref{tab:candidates} lists all 10
candidates with their predicted returns and profitability scores.
The remaining 17 tickers had negative CNN-LSTM predictions and were
correctly excluded from consideration.

\begin{table}[h]
\centering
\caption{Task~4.1 candidate stocks passing the profitability filter,
         sorted by profitability score descending.}
\label{tab:candidates}
\renewcommand{\arraystretch}{1.3}
\begin{tabular}{llrrrr}
\hline
\textbf{Ticker} & \textbf{Sector}
  & \textbf{Last Close} & \textbf{Pred.\ 7d}
  & \textbf{Exp.\ Return} & \textbf{Prof.\ Score} \\
\hline
STB & Banking    & 66,900 & 67,801 & +1.35\% & 0.0042 \\
FMC & Seafood    & 38,000 & 38,206 & +0.54\% & 0.0030 \\
RAL & Cables     & 87,600 & 88,021 & +0.48\% & 0.0027 \\
GMD & Logistics  & 72,600 & 72,871 & +0.37\% & 0.0018 \\
SFI & Logistics  & 28,450 & 28,543 & +0.33\% & 0.0013 \\
PET & Trading    & 48,500 & 48,692 & +0.40\% & 0.0012 \\
BMP & Materials  & 133,000 & 133,575 & +0.43\% & 0.0012 \\
PPC & Power      & 10,250 & 10,270 & +0.20\% & 0.0011 \\
DHA & Pharma     & 47,100 & 47,245 & +0.31\% & 0.0008 \\
DMC & Pharma     & 60,300 & 60,347 & +0.08\% & 0.0004 \\
\hline
\end{tabular}
\end{table}





\subsection{Task 4.2 --- Risk Scoring and Candidate Exclusion}

\subsubsection{Risk Metric Construction}

Five historical risk metrics are computed per ticker from the
full price history, each capturing a distinct dimension of downside
exposure:

\begin{enumerate}
  \item \textbf{Annualised volatility.} Rolling 252-day standard
    deviation of daily log-returns, scaled to annual. Captures
    overall price variability.
  \item \textbf{Maximum drawdown (MDD).} The largest peak-to-trough
    decline over the full history:
    \begin{equation}
      \text{MDD}_i = \max_{s < t} \frac{P_s - P_t}{P_s}
    \end{equation}
    MDD captures tail risk that volatility alone misses --- a stock
    can have moderate daily volatility but catastrophic multi-month
    drawdowns.
  \item \textbf{Downside deviation.} Semi-deviation computed only on
    negative daily returns (Sortino denominator), penalising downside
    moves disproportionately relative to symmetric volatility.
  \item \textbf{Value-at-Risk (VaR$_{5\%}$).} The 5th-percentile
    daily log-return, representing the expected loss on the worst
    5\% of trading days.
  \item \textbf{Conditional Value-at-Risk (CVaR$_{5\%}$).} The mean
    of all returns below VaR$_{5\%}$, capturing the average severity
    of extreme loss days.
\end{enumerate}

Each metric is normalised to $[0,1]$ via min-max scaling across the
27-ticker universe, and a composite risk score is computed as the
equal-weighted mean of the five normalised metrics:

\begin{equation}
\text{RiskScore}_i = \frac{1}{5}
  \sum_{m \in \{\sigma,\, \text{MDD},\, \text{DD},\, \text{VaR},\, \text{CVaR}\}}
  \widetilde{m}_i
\label{eq:risk_score}
\end{equation}

where $\widetilde{m}_i$ denotes the min-max normalised value of
metric $m$ for ticker $i$.

\subsubsection{High-Risk Exclusion}

Tickers with composite risk score $> 0.65$ are excluded from both
portfolio strategies. Table~\ref{tab:excluded} lists the five
excluded tickers. Notably, three of the five (PVD, TDH, VNE) are
cyclical or speculative stocks with histories of extreme drawdowns
exceeding 80\%, making them unsuitable for systematic allocation
regardless of short-term model signals.

\begin{table}[h]
\centering
\caption{High-risk tickers excluded from portfolio construction
         (composite risk score $> 0.65$).}
\label{tab:excluded}
\renewcommand{\arraystretch}{1.3}
\begin{tabular}{llc}
\hline
\textbf{Ticker} & \textbf{Sector} & \textbf{Risk Score} \\
\hline
PET & Petroleum Distribution & 0.864 \\
MHC & Real Estate            & 0.737 \\
VNE & Construction           & 0.722 \\
TDH & Real Estate            & 0.709 \\
PVD & Oil \& Gas             & 0.651 \\
\hline
\end{tabular}
\end{table}

After exclusion, the \textbf{surviving candidate universe} for
portfolio optimisation consists of the 10 profitability-filter
survivors minus any excluded by risk scoring. Of the original
10 candidates, none overlapped with the 5 high-risk exclusions,
leaving all 10 candidates available for MVO.

\subsection{Task 4.3 --- Mean-Variance Optimisation}

\textbf{Prudent portfolio} --- maximise Sharpe ratio, max single weight = 15\%:
\begin{equation}
\max_{\mathbf{w}} \; \frac{\mathbf{w}^\top \boldsymbol{\mu} - r_f}
                          {\sqrt{\mathbf{w}^\top \Sigma \mathbf{w}}}
\quad \text{s.t.}\ \sum_i w_i = 1,\; 0 \leq w_i \leq 0.15
\end{equation}

\textbf{Risk-taking portfolio} --- maximise expected return, Sharpe $\geq 0.80$
floor, max single weight = 30\%:
\begin{equation}
\max_{\mathbf{w}} \; \mathbf{w}^\top \boldsymbol{\mu}
\quad \text{s.t.}\ \frac{\mathbf{w}^\top \boldsymbol{\mu} - r_f}
                        {\sqrt{\mathbf{w}^\top \Sigma \mathbf{w}}} \geq 0.80,\;
\sum_i w_i = 1,\; 0 \leq w_i \leq 0.30
\label{eq:mvo_return}
\end{equation}


where $\boldsymbol{\mu}$ is the vector of expected returns (Task~2
7-day log-return predictions, annualised), $\Sigma$ is the sample
covariance matrix estimated from 252 trading days of historical
log-returns, and $r_f = 0.045$ is the risk-free rate (Vietnamese
1-year Treasury yield).




\subsubsection{Signal Tilt}

The raw MVO weights $\mathbf{w}^{\text{MVO}}$ are blended with a
signal conviction tilt derived from Task~3 BUY probabilities:

\begin{equation}
w_i^{\text{final}} = 0.70 \cdot w_i^{\text{MVO}}
+ 0.30 \cdot \frac{p_{\text{buy},i}}{\sum_j p_{\text{buy},j}}
\label{eq:signal_tilt}
\end{equation}




\subsection{Portfolio Results}

Table~\ref{tab:portfolios} compares the two portfolios.
The Prudent portfolio maximises the Sharpe ratio subject to a
15\% maximum single-weight constraint and a risk score ceiling
of 0.50 (Low--Medium tickers only). The Risk-Taking portfolio
relaxes both constraints (30\% max weight, risk ceiling 0.65)
to allow higher concentration in the highest-return candidates.
\begin{figure}[h]
    \centering
    \includegraphics[width=0.85\textwidth]{task4_portfolio_optimization.png}
    \caption{Efficient frontier and final portfolio allocations for both strategies.}
    \label{fig:t4_portfolio}
\end{figure}
\begin{figure}[h]
    \centering
    \includegraphics[width=0.85\textwidth]{projectedReturns.png}
    \caption{Projected 7-day cumulative returns for both portfolios
             using Task~2 consecutive-day predictions.}
    \label{fig:t4_projected_returns}
\end{figure}
\begin{table}[h]
\centering
\caption{Portfolio construction parameters and performance metrics.}
\label{tab:portfolios}
\renewcommand{\arraystretch}{1.3}
\begin{tabular}{lcc}
\hline
\textbf{Parameter} & \textbf{Prudent} & \textbf{Risk-Taking} \\
\hline
Objective           & Max Sharpe ratio & Max return (Sharpe $\geq 0.80$) \\
Max single weight   & 15\%             & 30\% \\
Risk score ceiling  & $\leq 0.50$      & $\leq 0.65$ \\
\hline
\textbf{Expected Annual Return} & \textbf{11.61\%} & \textbf{13.56\%} \\
\textbf{Annual Volatility}      & \textbf{13.67\%} & \textbf{15.88\%} \\
\textbf{Sharpe Ratio}           & \textbf{0.520}   & \textbf{0.571}   \\
Number of stocks    & 6                & 5 \\
\hline
\end{tabular}
\end{table}

Table~\ref{tab:portfolio_weights} reports the final weights for
both portfolios. Both select from the same low-to-medium risk
pool (FMC, RAL, PPC, DMC appear in both), confirming that the
optimiser consistently favours these diversified, liquid tickers.
The Risk-Taking portfolio concentrates more heavily in FMC, RAL,
and PPC (each $\approx$27\%) at the cost of dropping DHA and SFI.

\begin{table}[h]
\centering
\caption{Final signal-boosted portfolio weights.}
\label{tab:portfolio_weights}
\renewcommand{\arraystretch}{1.3}
\begin{tabular}{lcc}
\hline
\textbf{Ticker} & \textbf{Prudent} & \textbf{Risk-Taking} \\
\hline
FMC & 17.5\% & 27.2\% \\
RAL & 17.5\% & 27.2\% \\
PPC & 17.5\% & 25.8\% \\
DMC & 17.2\% & 13.0\% \\
SFI & 15.9\% & ---    \\
DHA & 14.4\% & ---    \\
GMD & ---    &  6.7\% \\
\hline
\end{tabular}
\end{table}

\subsection{Task 4.4 --- Target-Return Portfolio}

An interactive optimisation cell allows an investor to specify a
desired annual return $r_{\text{target}}$ and receive the
minimum-variance portfolio that exactly meets it:

\begin{equation}
\min_{\mathbf{w}}\; \mathbf{w}^\top \boldsymbol{\Sigma}\, \mathbf{w}
\quad\text{subject to}\quad
\mathbf{w}^\top \boldsymbol{\mu} = r_{\text{target}},\;\;
\sum_i w_i = 1,\;\;
0 \leq w_i \leq 40\%
\label{eq:min_var}
\end{equation}

The feasible return range for the current candidate universe is
\textbf{9.1\%} (Global Minimum Variance Portfolio) to \textbf{19.5\%}
(maximum single-stock return). The slider is bounded to
$[\max(5\%, r_{\min}),\; r_{\max}]$, preventing infeasible
sub-risk-free-rate targets. This feature operationalises the efficient
frontier concept: rather than presenting an abstract curve, it
returns a concrete, investable allocation for any valid return target.

\subsection{Methodology Justification}

Table~\ref{tab:t4_justification} summarises the key methodological
choices and their rationale.

\begin{table}[h]
\centering
\caption{Task~4 design choices and justification.}
\label{tab:t4_justification}
\renewcommand{\arraystretch}{1.4}
\begin{tabular}{p{4.5cm}p{8cm}}
\hline
\textbf{Choice} & \textbf{Reasoning} \\
\hline
MVO via \texttt{scipy} (no PyPortfolioOpt)
  & Avoids external dependencies; the optimisation is fully
    transparent and auditable; standard for academic submissions. \\
7-day horizon for $\boldsymbol{\mu}$
  & Consistent with Task~3 signal horizon; short enough that
    predictions remain actionable; annualised with factor 252/7. \\
30\% Task~3 signal tilt
  & Directly integrates the classification output; 70\% MVO
    weight preserves diversification properties. \\
252-day covariance window
  & Standard 1-year lookback; balances recency with statistical
    stability of the covariance estimate. \\
Sample covariance (252-day)
  & Standard 1-year lookback; balances recency with statistical
    stability; no shrinkage needed for a 5--10 stock universe. \\
Candidate filter at 50\% of $\tau^*$
  & Yields 8--12 candidates at any evaluation date, sufficient
    for diversification while excluding clearly negative signals. \\
High-risk exclusion before MVO
  & Ensures MVO cannot select a high-volatility stock simply
    because its predicted return is large. \\
Two portfolio variants
  & Explicitly addresses the risk-appetite dimension:
    Prudent for conservative investors, Risk-Taking for
    growth-oriented investors. \\
\hline
\end{tabular}
\end{table}

\subsection{Portfolio Results}

Table~\ref{tab:portfolio_alloc} presents the final weights for both strategies.

\begin{table}[h]
\centering
\caption{Final portfolio allocations after MVO and signal tilt.
         Weights are reported as percentage of total portfolio.}
\label{tab:portfolio_alloc}
\renewcommand{\arraystretch}{1.3}
\begin{tabular}{llrr}
\hline
\textbf{Ticker} & \textbf{Sector}
  & \textbf{Prudent (\%)} & \textbf{Risk-Taking (\%)} \\
\hline
FMC & Seafood    & 17.5 & 27.2 \\
RAL & Cables     & 17.5 & 27.2 \\
PPC & Utilities  & 17.5 & 25.8 \\
DMC & Pharma     & 17.2 & 13.0 \\
SFI & Logistics  & 15.9 & --- \\
DHA & Pharma     & 14.4 & --- \\
GMD & Logistics  & --- & 6.7 \\
STB & Banking    & --- & --- \\
BMP & Plastics   & --- & --- \\
ACB & Banking    & --- & --- \\
\hline
\textbf{Total}  & & \textbf{100\%} & \textbf{100\%} \\
\hline
\end{tabular}
\end{table}

The Prudent portfolio concentrates in six tickers with
comparable weights (14--18\%), reflecting the minimum-variance
objective's preference for diversification. The Risk-Taking
portfolio concentrates 80\% in three tickers (FMC, RAL, PPC),
reflecting the Sharpe-maximisation objective's tolerance for
concentration when predicted return-to-risk ratios are high.
The four tickers receiving zero weight in both strategies
(STB, BMP, ACB) had the lowest profitability scores among
the 10 candidates and were driven to zero by the MVO
long-only constraint.


\clearpage
\section{Task 5 --- Production Deployment}

\subsection{System Overview}

Task~5 packages all model outputs into a user-facing platform:
a FastAPI REST backend serving predictions, trading signals, and
portfolio data, paired with a React~19 dashboard and an automated
daily data pipeline. The full stack is containerised via Docker
Compose for reproducible deployment.

\subsection{Backend Architecture}

The FastAPI application (\texttt{api/main.py}) exposes all model
outputs under a versioned prefix \texttt{/api/v1}.
Table~\ref{tab:api_endpoints} summarises the primary endpoints.

\begin{table}[h]
\centering
\caption{Core API endpoints.}
\label{tab:api_endpoints}
\renewcommand{\arraystretch}{1.3}
\small
\begin{tabular}{llp{5.5cm}}
\hline
\textbf{Method} & \textbf{Endpoint} & \textbf{Description} \\
\hline
GET & \texttt{/predict/\{ticker\}?task=task2\_1} & VN CNN-LSTM next-day prediction \\
GET & \texttt{/predict/nasdaq/\{ticker\}?k=1}    & NASDAQ LSTM prediction \\
GET & \texttt{/signals/\{ticker\}}               & BUY / SELL signal + confidence \\
GET & \texttt{/signals/all/latest}               & Signals for all 27 VN tickers \\
GET & \texttt{/portfolio/summary}                & Both portfolio allocations \\
GET & \texttt{/market/history/\{ticker\}}        & OHLCV history via yfinance \\
GET & \texttt{/live/prices}                      & SSE real-time price stream \\
POST & \texttt{/pipeline/trigger}               & Manual pipeline run (JWT required) \\
POST & \texttt{/auth/register}                  & User registration \\
POST & \texttt{/auth/login}                     & Login, returns JWT \\
\hline
\end{tabular}
\end{table}

\textbf{Cache-first serving.}
All prediction endpoints follow a cache-first strategy: the API
first queries MongoDB for a cached prediction from the current
trading day; only if no cache entry exists does it run live
inference. This design ensures sub-millisecond response times for
cached predictions while preserving the ability to serve fresh
results on demand.

\textbf{In-memory model cache.}
All \texttt{.keras} model files and \texttt{.pkl} scalers are loaded
into memory at application startup via \texttt{preload\_all()} in the
FastAPI lifespan handler. Subsequent inference requests require no
disk I/O, reducing per-request latency to the GPU/CPU compute time
only.

\subsection{Automated Data Pipeline}

\subsubsection{APScheduler Daily Pipeline}

An \texttt{AsyncIOScheduler} (APScheduler 3.10) is started inside the
FastAPI lifespan and fires the \texttt{run\_pipeline()} coroutine
daily at \textbf{06:00 ICT} (\texttt{Asia/Ho\_Chi\_Minh} timezone),
30 minutes after HOSE market data becomes available. The pipeline
executes five sequential stages:

\begin{enumerate}
  \item \textbf{Ingest.} Download the latest OHLCV bars for all 58
    tickers (27 VN + 31 NASDAQ) from Yahoo Finance via
    \texttt{yfinance}.
  \item \textbf{Feature engineering.} Apply \texttt{\_add\_vn\_features()}
    (24 features) or \texttt{\_add\_nasdaq\_features()} (20 features)
    to the refreshed data.
  \item \textbf{Inference.} Run all Task~2 and Task~1 models on the
    latest 60-day (NASDAQ) or 30-day (VN) window.
  \item \textbf{Signal classification.} Run all 54 Task~3 classifiers
    and apply calibrated thresholds.
  \item \textbf{MongoDB upsert.} Write prediction documents to the
    \texttt{predictions} collection with a date key, overwriting any
    existing entry for the current day.
\end{enumerate}

The complete pipeline for all 58 tickers completes in approximately
41 seconds. A \texttt{pipeline\_runs} collection records each
execution's start time, duration, status, and any per-ticker errors,
accessible via \texttt{GET /api/v1/pipeline/status}.

\subsubsection{Manual Trigger}

If the container was inactive at 06:00 (e.g.\ after a restart), a
stale prediction state can be corrected via the authenticated
pipeline trigger endpoint:

\begin{verbatim}
POST /api/v1/pipeline/trigger
\end{verbatim}

The endpoint runs \texttt{run\_pipeline()} as a FastAPI
\texttt{BackgroundTask}, returning a 202 Accepted response
immediately while the pipeline executes asynchronously.
JWT authentication is required to prevent unauthenticated
re-training triggers in public deployments.

\subsection{Infrastructure Stack and Design Journey}

\subsubsection{Docker Compose Stack}

The production stack comprises three services orchestrated by
Docker Compose (Figure~\ref{fig:stack}):

\begin{itemize}
  \item \textbf{api} --- \texttt{python:3.11-slim} image;
    FastAPI served by Uvicorn on port~8000; mounts
    \texttt{models/} read-only and \texttt{nasdaq-historical-data/}
    read-write (pipeline writes refreshed CSVs here).
  \item \textbf{ui} --- \texttt{nginx:alpine} image; serves the
    Vite-built React bundle and proxies \texttt{/api/} requests
    to the API container.
  \item \textbf{mongo} --- \texttt{mongo:7} with a named volume;
    present only in the local \texttt{docker-compose.local.yml}
    override; production targets MongoDB Atlas.
\end{itemize}

\begin{figure}[h]
\centering
\begin{verbatim}
  Internet
     |
  [ngrok tunnel]  (local demo)  /  [DNS + TLS]  (production)
     |
  [nginx:alpine  :80]
     |-- /          --> React SPA bundle
     |-- /api/      --> proxy → api:8000
     |
  [python:3.11-slim  :8000]   FastAPI + APScheduler
     |-- /api/v1/predict      In-memory model cache
     |-- /api/v1/signals      yfinance ingestion
     |-- /api/v1/live         SSE stream
     |
  [mongo:7  :27017]   predictions · pipeline_runs · users
\end{verbatim}
\caption{Docker Compose service topology for the TensorFinance stack.}
\label{fig:stack}
\end{figure}

\subsubsection{Build Optimisation}

An initial Docker build context of \textbf{4.3~GB} caused build
times exceeding 8 minutes. Root cause analysis identified four
large directories inadvertently included in the build context:
\texttt{.venv/} (2.2~GB), \texttt{DlFinalProject/} (1.0~GB),
\texttt{ui/node\_modules/} (202~MB), and a \texttt{Docker.dmg}
installer (580~MB). Adding a \texttt{.dockerignore} file excluding
these paths reduced the build context to \textbf{$\approx$350~MB},
cutting build time to under 90 seconds.

\subsubsection{Stack Selection Rationale}

Three deployment configurations were evaluated before settling on
the final approach:

\textbf{Option 1 --- VPS (Hetzner CX22).}
A \$4.50/month virtual private server with Caddy TLS termination
and MongoDB Atlas would provide a permanent public URL without
dependency on a local machine staying awake. This option was
evaluated and a full deployment guide was prepared. However, the
remaining project timeline (one week) made VPS provisioning, DNS
configuration, and TLS certificate management impractical given
that the report had not yet been written.

\textbf{Option 2 --- Cloud notebook (Streamlit / Google Colab).}
Re-implementing the React dashboard in Streamlit would have
eliminated the Docker and Nginx complexity, but would have required
rewriting approximately 2,000 lines of frontend code. The existing
React dashboard was already functional and tested; the cost of
rewriting it was not justified.

\textbf{Option 3 --- Local Docker Compose + ngrok (chosen).}
The chosen approach runs all services locally and exposes port~80
to a public URL via ngrok's free tier. The demo URL
(\texttt{https://chooser-deprecate-briskly.ngrok-free.dev}) is
stable across container restarts as long as ngrok remains connected.
\texttt{caffeinate -i} prevents the host Mac from sleeping during
demos, ensuring container uptime. This configuration required zero
additional infrastructure and was fully operational within the
remaining timeline.

\subsection{Critical Bug: yfinance API Breakage}

During deployment testing, the data ingestion pipeline failed
for \textbf{all 31 NASDAQ tickers} with a \texttt{JSONDecodeError}.
Root cause: Yahoo Finance changed its API response format between
2024 and 2025, and \texttt{yfinance==0.2.48} (the pinned version)
had not been updated to handle the new format. Every ticker
returned an HTML error page instead of JSON, which the
\texttt{json.loads()} call could not parse.

The fix was a one-line change to \texttt{requirements.txt}:
\begin{verbatim}
yfinance>=0.2.54   # was: yfinance==0.2.48
\end{verbatim}

Version 1.3.0 resolved to by pip's constraint solver restored full
ingestion functionality. All NASDAQ CSVs were subsequently refreshed
manually before the Docker image rebuild.

\textbf{Lesson.} Hard-pinning third-party data provider clients
(yfinance, vnstock) introduces silent breakage risk when the
underlying API changes. The corrected \texttt{requirements.txt}
uses a minimum-version constraint rather than an exact pin for
these libraries.

\subsection{Authentication and Security}

The API implements JWT-based authentication using
\texttt{python-jose} and \texttt{bcrypt}. The \texttt{/auth/register}
and \texttt{/auth/login} endpoints issue tokens signed with a
configurable \texttt{JWT\_SECRET} environment variable. Only the
pipeline trigger endpoint requires authentication; all prediction,
signal, and portfolio endpoints are publicly readable to support
unauthenticated dashboard access.

Sensitive configuration (MongoDB connection string, JWT secret,
CORS origins) is injected exclusively via environment variables;
no secrets appear in the codebase or Docker image layers.

\section{Limitations and Future Work}

\subsection{Current Limitations}

\textbf{Directional accuracy below 50\% (Task~2).}
The CNN-LSTM is optimised for magnitude accuracy (Huber loss) rather
than directional accuracy. On days with low predicted returns, the
model correctly outputs a near-zero log-return, which minimises loss
regardless of sign. This results in a measured directional accuracy
of 44.4\% at $k=1$ --- below the naive 50\% baseline --- because
the binary direction metric penalises conservative magnitude-correct
predictions. A direction-aware loss (e.g.\ a weighted combination of
Huber and binary cross-entropy on the sign) could address this.

\textbf{No macro or news features.}
All features are derived from price and volume data only. Major price
moves driven by earnings surprises, central bank announcements, or
geopolitical events are not represented in the feature space. The model
treats these events as unexplained residuals, which inflates MAE on
high-volatility days and limits directional accuracy around event dates.

\textbf{Static portfolio weights.}
The Task~4 portfolio is computed once from a fixed evaluation-date
snapshot. In a live system, weights should be recomputed daily as
new CNN-LSTM predictions and Task~3 signals arrive. The current
implementation treats portfolio construction as a batch process
rather than a rolling one.

\textbf{No transaction costs or liquidity constraints.}
The MVO objective does not account for bid-ask spreads, market impact,
or minimum lot sizes. For low-liquidity VN stocks (e.g.\ PGC,
PET), the optimal theoretical allocation may not be executable at
the implied size.

\textbf{Public deployment availability.}
The ngrok free-tier URL resets whenever the local machine restarts
or the ngrok process is interrupted, making the public URL
non-persistent. A permanent deployment requires either a paid ngrok
subscription or a VPS with DNS.

\subsection{Future Work}

\textbf{Direction-aware training objective.}
Augmenting Huber loss with a directional penalty term
$\lambda \cdot \text{BCE}(\text{sign}(\hat{y}), \text{sign}(y))$
would directly optimise both magnitude and direction, addressing
the 44.4\% directional accuracy limitation at minimal architectural
cost.

\textbf{Macro feature integration.}
Incorporating macroeconomic features --- VN-Index return, USD/VND
exchange rate, world commodity prices --- as auxiliary inputs to the
LSTM could improve performance around systematic market moves. These
features are already stationary (log-returns and ratios) and would
slot directly into the existing feature pipeline.

\textbf{Rolling portfolio rebalancing.}
Extending the scheduler pipeline to recompute Task~4 weights daily
would transform the static portfolio snapshot into a live systematic
strategy, enabling backtesting with realistic rebalancing costs.

\textbf{VPS deployment.}
Migrating to a permanent VPS (Hetzner CX22 or equivalent) with
Caddy TLS termination and MongoDB Atlas would eliminate the ngrok
dependency and provide 24/7 availability. The Docker Compose
configuration is already production-ready; only DNS and TLS
provisioning would be required.

\textbf{Sentiment features.}
Integrating news sentiment scores (via a lightweight NLP classifier
on Vietnamese financial news) as additional features could capture
event-driven price moves that are invisible to pure OHLCV models.

\section{Conclusion}

This report has presented TensorFinance, an end-to-end deep learning
platform for financial time series forecasting and systematic trading,
developed across five progressive tasks in the CS313 curriculum.

Tasks~1 and~2 established the core modelling foundation. A three-layer
stacked LSTM trained on 31 NASDAQ tickers achieved a mean MAPE of 1.65\%
for next-day prediction, with consistent performance across bull and
bear market regimes. The CNN-LSTM hybrid developed for 27 Vietnamese HOSE
stocks achieved a mean MAPE of 1.31\%, demonstrating that the addition of
local CNN feature extraction and VN-specific candlestick features provides
meaningful advantage in a retail-dominated market. Both architectures share
a principled design foundation: all features are stationary by construction,
the target variable is a log-return rather than a raw price, and data leakage
is prevented via strict chronological splitting.

Task~3 extended the CNN-LSTM backbone to binary classification, generating
buy and sell trading signals per ticker. The key challenge --- a degenerate
all-positive collapse caused by regression weight transfer --- was resolved
through LSTM weight reinitialisation and focal loss, yielding mean AUC values
of 0.781 (BUY) and 0.806 (SELL) across 27 tickers.

Task~4 integrated prediction and signal outputs with classical Markowitz
mean-variance optimisation, producing two investable portfolio strategies
differentiated by risk objective. The
staged pipeline --- profitability filtering, risk exclusion, MVO, signal
tilt --- ensures that both return potential and downside risk are assessed
independently before weights are determined.

Task~5 packaged all model outputs into a production platform: a FastAPI backend
with daily APScheduler automation, a React~19 dashboard with live price streaming,
and a Docker Compose stack deployable on any Linux host. A critical data
provider breakage (yfinance API change) was diagnosed and resolved during
deployment, demonstrating the importance of minimum-version rather than exact
pinning for third-party API clients.

The project demonstrates that deep learning models can extract
predictive signal from financial time series when properly regularised and
evaluated under realistic conditions: strict temporal splits, scale-invariant
features, and robust loss functions. The directional accuracy limitation
identified in Task~2 and the static portfolio construction are the primary
targets for future work, alongside a permanent VPS deployment to replace
the current ngrok tunnelling setup.

















\end{document}